\documentclass[11pt]{article}
\usepackage[margin=1in]{geometry}
\usepackage{amsmath,amssymb,bm,mathtools}
\usepackage[T1]{fontenc}
\usepackage{lmodern}

\title{Detailed Note: Mixer Operator $U_B$ in the Computational Basis}
\author{}
\date{}

\begin{document}
	\maketitle
	
	\section*{1. Setup and Notation}
	For $n+1$ qubits, let computational-basis states be labeled by bit strings
	\[
	x,y \in \{0,1\}^{n+1}.
	\]
	Equivalently, in code these are indices in
	\[
	\mathrm{BasisIdx}(n) \cong \{0,1,\dots,2^{n+1}-1\}.
	\]
	
	A state is a function
	\[
	\psi : \{0,1\}^{n+1} \to \mathbb{C},
	\]
	with amplitude $\psi(y)$ on basis state $|y\rangle$.
	
	The depth-1 QAOA mixer generator is
	\[
	B = \sum_{j=1}^{n+1} X_j,
	\]
	so the mixer unitary is
	\[
	U_B(\beta) = e^{-i\beta B}.
	\]
	
	\section*{2. Kernel Representation}
	In the computational basis, $U_B(\beta)$ acts by a kernel (matrix element)
	\[
	K_\beta(x,y) := \langle x|U_B(\beta)|y\rangle.
	\]
	Hence for every output index $x$,
	\begin{equation}
	(U_B(\beta)\psi)(x) = \sum_{y} K_\beta(x,y)\,\psi(y).
	\label{eq:kernel-action}
	\end{equation}
	This is exactly the matrix-vector product in basis coordinates.
	
	\section*{3. Closed Form of the Kernel}
	Let $d(x,y)$ be the Hamming distance (number of bit positions where $x$ and $y$ differ). Then
	\begin{equation}
	K_\beta(x,y)
	= (\cos\beta)^{(n+1-d(x,y))}\,\bigl((-i)\sin\beta\bigr)^{d(x,y)}.
	\label{eq:kernel-closed-form}
	\end{equation}
	
	Interpretation:
	\begin{itemize}
		\item each equal bit contributes a factor $\cos\beta$,
		\item each flipped bit contributes a factor $(-i)\sin\beta$,
		\item only the number of differing bits matters, not their locations.
	\end{itemize}
	So $K_\beta$ depends on $(x,y)$ only through $d(x,y)$.
	
	\section*{4. Why This Is the Right Formula}
	Using commutativity of different-qubit Pauli operators,
	\[
	U_B(\beta) = \prod_{j=1}^{n+1} e^{-i\beta X_j}.
	\]
	For one qubit,
	\[
	e^{-i\beta X} = \cos\beta\,I - i\sin\beta\,X.
	\]
	In basis terms:
	\begin{itemize}
		\item no flip ($I$ channel): factor $\cos\beta$,
		\item flip ($X$ channel): factor $-i\sin\beta$.
	\end{itemize}
	Across $n+1$ qubits, to go from $|y\rangle$ to $|x\rangle$, exactly $d(x,y)$ flips are needed, giving \eqref{eq:kernel-closed-form}.
	
	\section*{5. Worked Example: Single Qubit ($n=0$)}
	Now $x,y\in\{0,1\}$ and
	\[
	K_\beta(x,y)=
	\begin{cases}
	\cos\beta, & x=y,\\
	-i\sin\beta, & x\neq y.
	\end{cases}
	\]
	So
	\[
	U_B(\beta)=
	\begin{pmatrix}
	\cos\beta & -i\sin\beta\\
	-i\sin\beta & \cos\beta
	\end{pmatrix}
	\]
	in the basis $(|0\rangle,|1\rangle)$.
	
	If $|s\rangle=|+\rangle=(|0\rangle+|1\rangle)/\sqrt2$, then
	\[
	U_B(\beta)|+\rangle
	= (\cos\beta - i\sin\beta)|+\rangle
	= e^{-i\beta}|+\rangle.
	\]
	So $|+\rangle$ is an eigenstate of the single-qubit mixer.
	
	\section*{6. Worked Example: Two Qubits ($n=1$)}
	Now $x,y\in\{00,01,10,11\}$ and
	\[
	K_\beta(x,y)=(\cos\beta)^{2-d(x,y)}((-i)\sin\beta)^{d(x,y)}.
	\]
	Hence:
	\begin{itemize}
		\item $d=0$: coefficient $(\cos\beta)^2$,
		\item $d=1$: coefficient $\cos\beta\,(-i\sin\beta)$,
		\item $d=2$: coefficient $(-i\sin\beta)^2$.
	\end{itemize}
	For fixed $x$, amplitudes are mixed from all four $y$ values according to these three distance classes.
	
	\section*{7. Connection to Your Lean Definitions}
	Your code encodes exactly this structure:
	\begin{itemize}
		\item \texttt{basisHammingDist n x y} = $d(x,y)$,
		\item \texttt{mixerEntry n beta x y} = $K_\beta(x,y)$ from \eqref{eq:kernel-closed-form},
		\item \texttt{U\_B n beta} applies \eqref{eq:kernel-action} by summing over $y$.
	\end{itemize}
	So yes: \textbf{$K_\beta$ is the mixer kernel}, i.e. the computational-basis matrix of $U_B(\beta)$.
	
\end{document}
